\section{Chương~\ref{sec:glutcomputer}. Nơ-ron \nt{GLUT} tính được gì}

Các khẳng định logic của chương là những định lý về mạch có trọng số không âm, được suy ra ngay trong chương; thực nghiệm xác nhận mô hình nơ-ron và việc các mạch được lắp theo những định lý đó (bộ phát hiện trùng khớp, đường trễ, nhóm tự duy trì, chuỗi) thực sự có trong não.

{\footnotesize
\setlength{\tabcolsep}{3pt}\renewcommand{\arraystretch}{1.15}
\begin{longtable}{>{\raggedright}p{0.5cm}>{\raggedright}p{4.0cm}>{\raggedright}p{5.6cm}>{\raggedright}p{2.3cm}>{\raggedright\arraybackslash}p{1.7cm}}
\toprule
TT & Khẳng định & Thí nghiệm và điều đã đo & Nguồn & Trạng thái \\
\midrule\endfirsthead
\toprule
TT & Khẳng định & Thí nghiệm và điều đã đo & Nguồn & Trạng thái \\
\midrule\endhead
1 & Nơ-ron là mạch RC có ngưỡng và đặt lại~\eqref{eq:glutneuron} & Bơm vào thân nơ-ron tháp vỏ não chuột cống (lát cắt) một dòng nhiễu giống dòng qua khớp thần kinh trong não sống. Sự phụ thuộc của tần số trung bình vào giá trị trung bình và độ phân tán của dòng được mô tả bằng nơ-ron tích phân rò nếu thêm sự thích nghi tần số chậm. Các khoảng giá trị của $\tau$ và độ trễ được nêu trong chương mà không dẫn thí nghiệm & \cite{Rauch2003} & đã đo \\
2 & Mô hình~\eqref{eq:glutneuron} là một sự đơn giản hóa: chính các nhánh đầu vào phát xung cục bộ & Ghi trực tiếp từ đuôi gai của nơ-ron tháp vỏ thị giác chuột nhắt in vivo: kích thích thị giác gây ra xung đuôi gai cục bộ phụ thuộc thụ thể NMDA; ức chế chúng làm nới rộng độ chọn lọc hướng của nơ-ron & \cite{Smith2013} & đã đo \\
3 & Cửa sổ trùng khớp $\Delta_{\max}=\tau\ln\frac{w}{\theta-w}$~\eqref{eq:window}; khi $\theta=1.5w$ nó bằng $0.7\tau$ & Hệ quả của mô hình~\eqref{eq:glutneuron}. Phép đo trực tiếp một cửa sổ cộng hẹp có ở nơ-ron có ức chế nhanh (chương~\ref{sec:gaba}, dòng~3 trong bảng của chương đó) & suy luận trong chương & lý thuyết \\
4 & Mạng chỉ gồm nơ-ron \nt{GLUT} chỉ tính được các hàm đơn điệu của đầu vào (mệnh đề~\ref{prop:monotone}) & Chứng minh trong chương: lặp từ trạng thái im lặng với vế phải không giảm. Đây là khẳng định về mô hình; chưa có thí nghiệm kiểm chứng nó trên mạng sống không có ức chế & suy luận trong chương & lý thuyết \\
5 & Giác quan cho ra một cặp dây: một số tế bào đáp ứng khi kích thích tăng, số khác khi kích thích giảm & Võng mạc: ngay từ tế bào lưỡng cực, tín hiệu của tế bào nón đã tách thành kênh bật và kênh tắt. Chẹn dược lý các tế bào lưỡng cực bật ở khỉ: các kênh đi riêng tới tận vỏ não; sau khi chẹn, khả năng phát hiện ánh sáng mạnh lên bị suy giảm, còn phát hiện ánh sáng yếu đi thì không & \cite{Schiller1992} & đã đo \\
6 & Với mã hai dây, mạng nơ-ron \nt{GLUT} tính mọi hàm logic một cách bất đồng bộ, không có nhịp chung, với cái giá là số nơ-ron gấp đôi (mệnh đề~\ref{prop:dualrail}, \eqref{eq:dualrail}, hình~\ref{fig:dualxor}) & Mã hai dây và việc xác định trạng thái sẵn sàng qua chính tín hiệu là kỹ thuật tiêu chuẩn của mạch bất đồng bộ không nhạy với độ trễ. Mạch XOR sáu nơ-ron được lắp trong chương. Chưa có thí nghiệm cho thấy não tính các hàm logic đúng bằng mã hai dây & \cite{Sparso2001} & lý thuyết \\
7 & Hiệu thời gian âm thanh đến hai tai lớn nhất ở người là $\approx0.6\ms$, còn não phân biệt được khoảng mười micro giây & Người nghe trong thí nghiệm hai phương án trả lời: ngưỡng phân biệt hiệu thời gian (75\,\% đúng) là $9\,\mu\text{s}$ với nhiễu dải $150$--$1700$\,Hz, $11\,\mu\text{s}$ với âm $1$\,kHz, $28\,\mu\text{s}$ với tiếng tách. Giới hạn $0.6\ms$ là phép tính của chương, $0.22\,\text{m}/343\,\text{m/s}$ & \cite{KlumppEady1956} & đã đo \\
8 & Ở chim, hiệu thời gian được biến thành vị trí nhờ đường trễ và bộ phát hiện trùng khớp~\eqref{eq:itd} (hình~\ref{fig:itd}) & Cú lợn lưng xám: sợi trục từ hai tai đi vào nhân của các bộ phát hiện trùng khớp từ hai phía ngược nhau; thời điểm xung đến dọc sợi trục thay đổi có trật tự theo độ sâu, và dải độ trễ qua nhân ($160\,\mu\text{s}$ trên $700\um$) trùng với dải độ trễ giữa hai tai & \cite{Carr1990} & đã đo \\
9 & Ở động vật có vú không tìm thấy hàng độ trễ; cùng bài toán được giải bằng bộ phát hiện trùng khớp với ức chế đúng thời điểm từ nơ-ron \nt{GLYC} & Ghi in vivo ở nhân trám trên giữa của chuột nhảy: đáp ứng không tạo thành bản đồ hướng; đưa glycine và chất chẹn nó qua vi pipet cho thấy sự ức chế \emph{glycine} đến đúng thời điểm quy định dải độ trễ làm việc. Ức chế ở đây là từ \nt{GLYC} chứ không phải từ \nt{GABA} & \cite{Brand2002,Grothe2010} & đã đo \\
10 & Nhóm có hồi tiếp mạnh là một flip-flop; hoạt động tự duy trì kéo dài hàng giây sau kích thích (hình~\ref{fig:bistable}) & Vỏ não trước trán của khỉ trong nhiệm vụ đưa mắt trì hoãn tới một điểm được ghi nhớ (trì hoãn $1$--$6$\,s): $87$ trên $170$ nơ-ron liên quan đến nhiệm vụ thay đổi tần số trong thời gian trì hoãn, ở $79\,\%$ trong số đó điều này phụ thuộc vào hướng tới điểm được ghi nhớ. Việc hoạt động này được duy trì bởi hồi tiếp có hai trạng thái bền là lý thuyết & \cite{Funahashi1989,Wang2001} & đã đo \\
11 & Bit được ghi nếu dòng vào kéo dài hơn $\approx0.7\tau$ (hình~\ref{fig:recurrentgroup}b) & Giải phương trình $\tau\,\dd r/\dd t=-r+f(wr+I)$ bằng phương pháp Euler với $w=1$, $I=0.6$; không có script trong \texttt{models/}. Chưa có thí nghiệm & tính toán trong chương & mô hình \\
12 & Bộ nhớ có thể lưu trong sự tạo thuận của khớp thần kinh, gần như không cần xung & Lý thuyết: calci dư trong đầu tận cùng giữ bản ghi khoảng một giây. Cơ sở: trong vỏ não trước trán giữa của chồn sương (ghi đồng thời nhiều nơ-ron) có một mạng con nơ-ron tháp nối bằng khớp thần kinh tạo thuận có hậu hiệu ứng kéo dài. Tổng quan các quan sát về những khoảng ``lặng'' khi giữ thông tin trong trí nhớ. Vỏ não dùng cách nào khi nào là câu hỏi còn mở & \cite{Mongillo2008,WangMarkram2006,Stokes2015} & giả thuyết \\
13 & Bộ nhớ bị xóa bởi sự mỏi: dự trữ chất dẫn truyền cạn kiệt & Ghi cặp nơ-ron tháp vỏ não: đáp ứng của khớp thần kinh giảm khi xung dày, tốc độ giảm được quy định bởi xác suất nhả. Chưa có bằng chứng trực tiếp rằng chính điều này dập tắt nhóm tự duy trì & \cite{Tsodyks1997} & đã đo \\
14 & Chuỗi các nhóm là bộ định thời khởi động theo sự kiện; ở chim biết hót, nơ-ron phát xung lần lượt nghiêm ngặt & Ghi các nơ-ron đã được nhận dạng ở trung tâm phát âm cao cấp của chim sẻ vằn khi ngủ và khi hót: mỗi nơ-ron gửi đầu ra tới nhân vận động phát một chùm xung ngắn vào một thời điểm chính xác của bài hót, và các nơ-ron phát xung lần lượt & \cite{Hahnloser2002} & đã đo \\
\bottomrule
\end{longtable}}
